\begingroup

\setlength{\LTcapwidth}{\textwidth}

\begin{longtable}{@{}>{\raggedright\arraybackslash}p{8.756cm}>{\raggedleft\arraybackslash}p{1.800cm}>{\raggedleft\arraybackslash}p{2.600cm}@{}}

\caption{Karakteristik responden evaluasi ahli.}

\label{tab:tab6.16_karakteristik_responden} \\

\toprule

\textbf{Karakteristik} & \textbf{Cacah} & \textbf{Persentase} \\

\midrule

\endfirsthead



\multicolumn{3}{@{}l}{\small\emph{Tabel \thetable{} (lanjutan)}} \\

\toprule

\textbf{Karakteristik} & \textbf{Cacah} & \textbf{Persentase} \\

\midrule

\endhead



\bottomrule

\endlastfoot



Dokter Spesialis & 3 & 25,0\% \\



Dokter & 1 & 8,3\% \\



Dokter Muda & 7 & 58,3\% \\



Mahasiswa kedokteran/kesehatan & 1 & 8,3\% \\



Pernah menggunakan CDSS/aplikasi kesehatan digital & 3 & 25,0\% \\



Belum pernah menggunakan CDSS/aplikasi kesehatan digital & 9 & 75,0\% \\



Pernah menangani pasien diabetes dengan CGM & 1 & 8,3\% \\



Jarang menangani pasien diabetes dengan CGM & 1 & 8,3\% \\



Belum pernah menangani pasien diabetes dengan CGM & 10 & 83,3\% \\

\end{longtable}



\noindent\tabnotestyle Bidang okupasi diisi 8 dari 12 responden: 3 Penyakit Dalam, 1 Dokter umum, 1 Kedokteran umum, 1 Koas, 1 Obgin, 1 Umum.\normalsize

\endgroup

